\begin{helper}
Fix an activation set $A$, a vertex $v$, and two levels $0\le a\le b\le 1$.
Let $\nu$ be a measure satisfying the rectangular priority clause from
\noderef{RandomIndependentSetSampling}.  Then the difference of the two
lower-rectangle values
\[
\nu[\omega_1=A]\,b\,a^m-\nu[\omega_1=A]\,a\,a^m,
\]
where
\[
m=\left|\{u\in V(G):u\in A\text{ and }u\sim_G v\}\right|.
\]
is at most the $\nu$-measure of the slice where the activation set is $A$,
the priority of $v$ lies in $(a,b]$, every activated neighbor of $v$ has
priority at most $a$, and all remaining priority coordinates lie in $[0,1]$.
\end{helper}
\begin{proof}
Let $R_b$ be the lower rectangle in which $\omega_1=A$, the coordinate of
$v$ is at most $b$, each activated neighbor of $v$ is at most $a$, and every
other coordinate is at most $1$.  Define $R_a$ in the same way, except that
the coordinate of $v$ is at most $a$.  Since $a\le b$, we have
$R_a\subseteq R_b$.

The stated slice is exactly $R_b\setminus R_a$.  Indeed, membership in
$R_b$ supplies the constraints with $\omega_2(v)\le b$; failing to belong to
$R_a$ is precisely the failure of $\omega_2(v)\le a$, because all other
coordinate constraints are identical in the two rectangles.  Hence the slice
condition is equivalent to $a<\omega_2(v)\le b$ together with the common
constraints.

By the rectangular computation \noderef{SamplingOneVertexPriorityRectangle},
\[
\nu(R_b)=\nu[\omega_1=A]\,
b\,a^m
\quad\text{and}\quad
\nu(R_a)=\nu[\omega_1=A]\,
a\,a^m.
\]
For arbitrary measures, the general measure inequality
\[
\nu(R_b)-\nu(R_a)\le \nu(R_b\setminus R_a)
\]
holds.  Substituting the two displayed rectangle values and the identification
of the slice with $R_b\setminus R_a$ proves the claimed lower bound for the
slice measure.
\end{proof}
